% einstein_photoelectric_ML_slides_v2.tex
% Beamer presentation: "Einstein Was Doing Machine Learning"
% Thomas J. Sargent, 2026
% Style: default theme (matching AI_past_present_future_PHBS_2025.tex)

\documentclass[aspectratio=169]{beamer}

\usetheme{default}
\setbeamertemplate{navigation symbols}{}

\usepackage{amsmath,amssymb}
\usepackage{booktabs}
\usepackage{array}
\usepackage{multirow}
\usepackage{microtype}
\usepackage{hyperref}
\usepackage{graphicx}

\hypersetup{colorlinks=true, linkcolor=blue, citecolor=blue, urlcolor=blue}

\title{Einstein Was Doing Machine Learning \\
{\normalsize The Photoelectric Effect as Statistical Inference}}
\author{Thomas J.\ Sargent \\
New York University \& Hoover Institution}
\date{2026}

%----------------------------------------------------------------------
\begin{document}
%----------------------------------------------------------------------

\begin{frame}
  \titlepage
\end{frame}

%----------------------------------------------------------------------
\begin{frame}{The Central Claim}

\begin{block}{Thesis}
When Einstein proposed the light-quantum hypothesis in 1905 to explain
the photoelectric effect, he was---in the language we now use---doing
\textbf{machine learning} and \textbf{statistical inference}.
\end{block}

\bigskip
Three acts:
\begin{enumerate}
  \item \textbf{Theory} (1905): Einstein proposes the model that defines
    the hypothesis class
  \item \textbf{Data quality} (1887--1916): experimenters spend 29 years
    collecting clean training data
  \item \textbf{Estimation} (1916): Millikan runs the regression and
    recovers Planck's constant to 0.5\%
\end{enumerate}

\end{frame}

%----------------------------------------------------------------------
\begin{frame}{Background: The Photoelectric Effect}

The phenomenon (Hertz, 1887; Hallwachs, 1888):
\begin{itemize}
  \item Shine light on a metal surface; electrons are ejected
  \item Apply a \emph{stopping potential} $V_0$ to just halt the
    fastest electrons
\end{itemize}

\medskip
The puzzle (Lenard, 1902):
\begin{itemize}
  \item $V_0$ rises with light \emph{frequency} $\nu$
  \item $V_0$ is \textbf{independent} of light \emph{intensity} $I$
  \item Classical wave theory predicted the opposite
\end{itemize}

\medskip
Classical prediction (wrong):
\begin{itemize}
  \item Wave amplitude $\propto \sqrt{I}$ delivers energy
    $\Rightarrow$ $V_0$ should rise with $I$
  \item $V_0$ should be independent of $\nu$
\end{itemize}

\end{frame}

%----------------------------------------------------------------------
\begin{frame}{Lenard's Finding as Feature Engineering}

Lenard ran an ablation experiment: measured $V_0$ as a function of
both $\nu$ and $I$.

\medskip
\begin{center}
\begin{tabular}{lcc}
\toprule
Feature & Effect on $V_0$ & Importance \\
\midrule
Frequency $\nu$ & Strong $\uparrow$ & \textbf{High} \\
Intensity $I$  & None             & \textbf{Zero} \\
\bottomrule
\end{tabular}
\end{center}

\medskip
Conclusion: drop $I$ from the feature set; keep $\nu$.

\bigskip
In modern ML language:
\begin{itemize}
  \item \textbf{Input features:} $\nu$, $I$, $\lambda$, \ldots
  \item \textbf{Label:} stopping potential $V_0$
  \item \textbf{Result:} feature-importance analysis eliminates $I$;
    selects $\nu$ as the sole relevant input
\end{itemize}

\end{frame}

%----------------------------------------------------------------------
\begin{frame}{Einstein's 1905 Insight: The Model Architecture}

The quantum hypothesis: each light quantum carries energy $E = h\nu$.

\medskip
A single quantum is absorbed by a single electron.
If $h\nu > W$ (work function), the electron escapes.

\medskip
Energy conservation gives:
\[
  eV_0 = h\nu - W
\]

\medskip
\begin{block}{The Prediction}
\[
  \boxed{V_0 = \frac{h}{e}\,\nu - \frac{W}{e}}
\]
\begin{itemize}
  \item Linear in frequency
  \item Slope $= h/e$ is \textbf{universal} across all metals
  \item Written before a single reliable $V_0$-vs-$\nu$ measurement existed
\end{itemize}
\end{block}

\end{frame}

%----------------------------------------------------------------------
\begin{frame}{The Inductive Bias}

The quantum hypothesis restricts the hypothesis class from \emph{all
functions} of $(\nu, I, \lambda, \ldots)$ to a single affine
function in $\nu$ alone, with a specific universal slope.

\bigskip
\begin{center}
\begin{tabular}{lll}
\toprule
Hypothesis class & Size & Data needed \\
\midrule
All measurable functions of $(\nu, I, \ldots)$ & Infinite-dimensional & $\gg 10^6$ \\
Affine functions of $\nu$ alone (Einstein) & 2-dimensional & $\sim 5$ \\
\bottomrule
\end{tabular}
\end{center}

\bigskip
The key intellectual act was not fitting the line---it was
knowing which line to fit.

\end{frame}

%----------------------------------------------------------------------
\begin{frame}{The Regression Equation}

Adding measurement noise and a nuisance parameter:
\[
  V_{0i} = \underbrace{\alpha}_{\text{intercept}} +
            \underbrace{\beta}_{h/e}\,\nu_i +
            \underbrace{\varepsilon_i}_{\text{noise}},
  \qquad
  \alpha = -\frac{W}{e} + \underbrace{v_c}_{\text{contact EMF}}
\]

\medskip
Parameters of interest:
\begin{itemize}
  \item $\beta = h/e$ --- slope $\Rightarrow$ Planck's $h$
  \item $\alpha$ --- intercept (metal-specific)
\end{itemize}

\medskip
Nuisance parameter: $v_c$, the contact EMF between emitter and collector.
\begin{itemize}
  \item Can be as large as 3\,V---comparable to $V_0$ itself
  \item Shifts intercept; \textbf{does not bias slope} if constant within a run
\end{itemize}

\end{frame}

%----------------------------------------------------------------------
\begin{frame}{The Supervised Learning Problem}

\begin{itemize}
  \item \textbf{Training data:} $\{(\nu_i, V_{0i})\}_{i=1}^n$
  \item \textbf{Task:} estimate $\beta = h/e$
  \item \textbf{Loss:} squared error
  \item \textbf{Hypothesis class:} affine functions of $\nu$
  \item \textbf{Inductive bias:} $\beta > 0$; universal across metals
\end{itemize}

\bigskip
The regression is trivial once you have clean data.

\bigskip
\textbf{The 29-year delay reflects data engineering, not statistical
sophistication.}

\end{frame}

%----------------------------------------------------------------------
\begin{frame}{29 Years of Collecting Training Data (1887--1916)}

\begin{center}
\small
\begin{tabular}{p{3.2cm}cp{3.5cm}p{4.3cm}}
\toprule
Experimenter & Year & Data quality & Problem \\
\midrule
Hertz / Hallwachs & 1887--88 & Qualitative only & No labeled $(\nu,V_0)$ pairs \\
Lenard            & 1902     & Qualitative       & No reliable numerical $V_0$ \\
Hughes            & 1912     & 3 UV lines        & Contact EMF uncontrolled; 65\% scatter in $h$ \\
Richardson \& Compton & 1912 & 8 metals          & Same problems; 60\% scatter \\
\midrule
\textbf{Millikan} & \textbf{1916} &
  \textbf{Clean surfaces, wide $\Delta\nu$, controlled $v_c$} &
  \textbf{$h$ to 0.5\%} \\
\bottomrule
\end{tabular}
\end{center}

\end{frame}

%----------------------------------------------------------------------
\begin{frame}{Millikan's Experimental Design (1916)}

Key data-quality innovations (``a machine shop in vacuo''):

\begin{itemize}
  \item \textbf{Fresh surfaces:} electromagnetic knife shaved alkali
    metal in vacuum---no oxidation, no work-function contamination
  \item \textbf{Spectrally pure lines:} mercury arc with filters to
    exclude UV companions
  \item \textbf{Wide frequency range:} 5461\,\AA\ to 2535\,\AA\
    ($\Delta\nu = 6.3\times10^{14}$\,Hz, nearly a factor of two)
  \item \textbf{Controlled nuisance:} Kelvin double-balance measured
    $v_c$ continuously; stable to $\pm 0.01$\,V over 9 hours
\end{itemize}

\medskip
Why the wide frequency range matters (Cram\'er--Rao bound):
\[
  \mathrm{SE}(\hat\beta) = \frac{\sigma}{\sqrt{S_{\nu\nu}}},
  \quad S_{\nu\nu} = \sum_i(\nu_i - \bar\nu)^2
\]
Wider $\Delta\nu$ $\Rightarrow$ larger $S_{\nu\nu}$ $\Rightarrow$ smaller SE.

\end{frame}

%----------------------------------------------------------------------
\begin{frame}{Millikan's Training Data}

\begin{center}
\begin{tabular}{rrcc}
\toprule
$\lambda$ (\AA) & $\nu$ ($10^{14}$\,Hz)
  & $V_0^{\,\mathrm{Na}}$ (V) & $V_0^{\,\mathrm{Li}}$ (V) \\
\midrule
5461.0 &  5.490 & 0.454  & ---   \\
4339.4 &  6.908 & 1.039  & 0.499 \\
4046.8 &  7.409 & 1.246  & 0.706 \\
3650.2 &  8.213 & 1.578  & 1.037 \\
3125.5 &  9.591 & 2.147  & 1.606 \\
2534.7 & 11.827 & 3.030  & 2.529 \\
\bottomrule
\end{tabular}
\end{center}

\medskip
True stopping potential $= V_{0,\text{obs}} + v_c$
\hspace{1em} (Na: $v_c = 2.51$\,V;\; Li: $v_c = 1.52$\,V)

\medskip
\begin{itemize}
  \item 6 labeled training examples per metal; two metals
  \item Points span nearly a factor of two in frequency
\end{itemize}

\end{frame}

%----------------------------------------------------------------------
\begin{frame}{OLS Regression: Estimating Planck's Constant}

Primary regression: sodium, five points

\medskip
Standard OLS formula:
\[
  \hat\beta = \frac{S_{\nu V}}{S_{\nu\nu}}
            = \frac{3.830}{9.277}
            = \mathbf{4.129\times10^{-15}}\,\text{V$\cdot$s}
\]

\medskip
\begin{center}
\begin{tabular}{p{4.5cm}rr}
\toprule
 & Value & \\
\midrule
$R^2$ & $\approx 0.99999$ & \\
SE$(\hat\beta)$ & $4.6\times10^{-18}$\,V$\cdot$s & \\
Implied $h$ (Millikan's $e$, 1916) & $6.573\times10^{-34}$\,J$\cdot$s & $-0.80\%$ \\
Implied $h$ (modern NIST $e$) & $6.614\times10^{-34}$\,J$\cdot$s & $-0.18\%$ \\
Modern exact $h$ & $6.626\times10^{-34}$\,J$\cdot$s & \\
\bottomrule
\end{tabular}
\end{center}

\medskip
$h$ recovered to \textbf{0.5\%} of its modern value from \textbf{5 data points}.

\end{frame}

%----------------------------------------------------------------------
\begin{frame}{MLE $=$ OLS Under Gaussian Errors}

Parametric model:
\[
  V_{0i} = \alpha + \beta\,\nu_i + \varepsilon_i,
  \quad
  \varepsilon_i \overset{\text{iid}}{\sim} \mathcal{N}(0,\sigma^2)
\]

Log-likelihood:
\[
  \ell(\alpha,\beta,\sigma) =
  -\frac{n}{2}\ln(2\pi) - n\ln\sigma
  - \frac{1}{2\sigma^2}\sum_{i}(V_{0i}-\alpha-\beta\nu_i)^2
\]

Maximising $\ell$ over $(\alpha,\beta)$ minimises RSS---exactly OLS:
\[
  \boxed{(\hat\alpha_{\mathrm{MLE}},\hat\beta_{\mathrm{MLE}}) =
         (\hat\alpha_{\mathrm{OLS}},\hat\beta_{\mathrm{OLS}})}
\]

\medskip
The MSE training loss used in neural-network regression \emph{is} the
Gaussian log-likelihood---the physics and the ML objective coincide.

\end{frame}

%----------------------------------------------------------------------
\begin{frame}{The Cram\'er--Rao Bound and Experimental Design}

Fisher information and Cram\'er--Rao bound for $\beta = h/e$:
\[
  \mathcal{I}_{\beta\beta} = \frac{S_{\nu\nu}}{\sigma^2},
  \qquad
  \mathrm{Var}(\tilde\beta) \geq \frac{\sigma^2}{S_{\nu\nu}}
\]
OLS achieves this bound (Gauss--Markov + Cram\'er--Rao).

\medskip
\begin{center}
\small
\begin{tabular}{p{5.3cm}rrr}
\toprule
Dataset & $n$ & $S_{\nu\nu}$ ($10^{28}$ Hz$^2$) & SE($\hat\beta$) ($10^{-18}$\,V$\cdot$s) \\
\midrule
Restricted range (like Hughes) & 3 & 1.96  & 10.0 \\
Millikan 5-pt                  & 5 & 9.28  &  4.6 \\
Millikan full 6-pt             & 6 & 28.14 &  2.6 \\
\bottomrule
\end{tabular}
\end{center}

\medskip
Design lesson: span frequencies widely.

Hughes's narrow range gives a standard error \textbf{2.2$\times$ larger}
than Millikan's full range.

\end{frame}

%----------------------------------------------------------------------
\begin{frame}{Universality: Zero-Shot Transfer Across Metals}

The quantum hypothesis predicts $\beta = h/e$ is the same for
\emph{all} metals---a strong, falsifiable structural restriction.

\medskip
Separate regressions:
\begin{align*}
  \hat\beta_{\mathrm{Na}} &= 4.129\times10^{-15}\;\text{V$\cdot$s} \\
  \hat\beta_{\mathrm{Li}} &= 4.137\times10^{-15}\;\text{V$\cdot$s} \\
  \text{difference} &= 0.19\%
\end{align*}

Wald test of $H_0\colon \beta_{\mathrm{Na}} = \beta_{\mathrm{Li}}$:\quad
$z \approx 0.87$ ($\ll 1.96$). Cannot reject universality.

\medskip
In ML language:
\begin{itemize}
  \item \textbf{Pre-training:} fit slope on sodium
  \item \textbf{Zero-shot transfer:} apply slope to lithium
  \item \textbf{Fine-tuning needed:} only the intercept (work function)
  \item \textbf{Transfer error:} 0.2\%
\end{itemize}

\end{frame}

%----------------------------------------------------------------------
\begin{frame}{Historical Estimates of Planck's Constant}

\begin{center}
\begin{tabular}{p{4.4cm}rcr}
\toprule
Authors & Year & $h$ ($10^{-27}$\,erg$\cdot$s) & Error \\
\midrule
Hughes              & 1912 & 3.55--5.85        & $-12$ to $-46\%$ \\
Richardson \& Compton & 1912 & 3.5--5.8        & $-12$ to $-47\%$ \\
Planck (black body) & 1913 & 6.415             & $-3.2\%$         \\
\midrule
Millikan (Na, 5-pt) & 1916 & 6.573             & $-0.80\%$        \\
Millikan (Na, 9-pt) & 1916 & 6.578             & $-0.73\%$        \\
Millikan (Li)       & 1916 & 6.584             & $-0.64\%$        \\
\textbf{Millikan mean} & \textbf{1916} & \textbf{6.579} & $\mathbf{-0.72\%}$ \\
\midrule
Modern NIST 2022    & 2022 & 6.6261            & 0                \\
\bottomrule
\end{tabular}
\end{center}

\medskip
Pre-Millikan estimates: biased downward 10--47\%.

Three causes: stray UV, narrow $\Delta\nu$, uncontrolled $v_c$.

Millikan solved all three.

\end{frame}

%----------------------------------------------------------------------
\begin{frame}{Millikan's Epistemic Reluctance}

Millikan confirmed Einstein's equation to 0.5\%. Then wrote:

\bigskip
\begin{quote}
``Despite then the apparently complete success of the Einstein equation,
the physical theory of which it was designed to be the symbolic expression
is found \textbf{so untenable} that Einstein himself, I believe, no longer
holds to it.''
\end{quote}

\quad \quad \quad Millikan (1916), p.\ 388

\bigskip
He had hoped the experiment would \textbf{refute} the equation.

\bigskip
\begin{itemize}
  \item Confirming a regression equation is \textbf{not} the same as
    confirming the structural model that generated it
  \item Compton scattering (1923) finally convinced him
  \item Nobel acceptance (1923): ``\ldots first direct experimental proof
    \ldots of the complete validity of the Einstein equation''
\end{itemize}

\end{frame}

%----------------------------------------------------------------------
\begin{frame}{The ML-to-Physics Dictionary}

\begin{center}
\small
\begin{tabular}{p{5.0cm}p{4.8cm}p{2.6cm}}
\toprule
\textbf{Scientific step} & \textbf{ML name} & \textbf{Who} \\
\midrule
Propose $E = h\nu$; derive $V_0 = (h/e)\nu - W/e$
  & Inductive bias + hypothesis class & Einstein (1905) \\
$V_0$ depends on $\nu$ not $I$
  & Feature engineering & Lenard (1902) \\
Measure $(\nu_i, V_{0i})$ pairs
  & Labeled training data & Millikan (1916) \\
Control contact EMF $v_c$
  & Nuisance-variable identification & Millikan (1916) \\
Min.\ $\sum(V_{0i} - \hat V_{0i})^2$
  & MSE training (OLS/MLE) & Millikan (1916) \\
$|\hat\beta - h/e| < 1\%$
  & Test accuracy & Millikan (1916) \\
Na and Li give same $\hat\beta$
  & Zero-shot transfer learning & Millikan (1916) \\
\bottomrule
\end{tabular}
\end{center}

\end{frame}

%----------------------------------------------------------------------
\begin{frame}{The Hypothesis Class}

Einstein's hypothesis class:
\[
  \mathcal{H} = \bigl\{f_{\alpha,\beta}: \nu \mapsto \alpha + \beta\nu
  \bigm| \alpha\in\mathbb{R},\;\beta>0\bigr\}
\]
with the cross-metal restriction that $\beta$ is universal.

\medskip
Properties:
\begin{itemize}
  \item VC dimension $= 2$
  \item One-parameter family of parallel lines
  \item Derived \emph{deductively} from $E = h\nu$
  \item Five training points are more than sufficient
\end{itemize}

\medskip
What a naive ML system misses:
\begin{enumerate}
  \item Cannot identify $\hat\beta = 4.129\times10^{-15}$ as $h/e$
  \item Cannot impose $\beta_{\mathrm{Na}} = \beta_{\mathrm{Li}}$
    from first principles
  \item Cannot decompose intercept into $W/e$ vs.\ $v_c$
  \item Cannot reduce data requirement from $O(10^3)$ to 5
\end{enumerate}

\end{frame}

%----------------------------------------------------------------------
\begin{frame}{Einstein vs.\ a Deep Neural Network}

\begin{center}
\small
\begin{tabular}{p{2.5cm}rp{2.7cm}p{2.8cm}p{3.5cm}}
\toprule
Epoch & $n$ & Einstein model & DNN (VC $\approx$ 5700) & DNN outcome \\
\midrule
Lenard 1902   & $\approx 0$  & Sufficient$^*$ & Insufficient
  & Cannot train \\
Hughes 1912   & $\sim 30$    & Sufficient
  & Insufficient ($\times 190$) & Memorises; no gen. \\
Millikan 1916 & 5            & Sufficient
  & Insuf.\ ($\times 1140$) & Interpolates only \\
Modern era    & $\sim 5000$  & Sufficient
  & Sufficient & Recovers line; cannot identify $h/e$ \\
\bottomrule
\end{tabular}
\end{center}

\medskip
{\small $^*$Functional form from theory; only two free parameters remain.}

\medskip
A three-layer DNN (16 neurons/layer): $p \approx 593$ parameters,
VC dim $\approx 5700$. Reliable 1\% generalisation requires
$n \gtrsim 5700$---short by three orders of magnitude in 1916.

\end{frame}

%----------------------------------------------------------------------
\begin{frame}{The DNN on Millikan's Five Points}

What happens if you train any DNN on $n = 5$ data points?

\medskip
Interpolation regime (any DNN with $p > 5$):
\begin{itemize}
  \item Training loss $= 0$ (memorises all points)
  \item Function is unconstrained between points
  \item Cannot estimate slope uncertainty: zero residual df
\end{itemize}

\medskip
OLS with $n=5$, $p=2$:
\begin{itemize}
  \item 3 residual degrees of freedom
  \item $\mathrm{SE}(\hat\beta) = 4.6\times10^{-18}$\,V$\cdot$s
  \item 0.11\% relative precision, at the Cram\'er--Rao bound
\end{itemize}

\medskip
Einstein's theory---not more data, not a better algorithm---is what made
five training points sufficient.

\medskip
DNNs win for problems without a known parsimonious parametric family
(images, proteins, language). The photoelectric problem sits at the
\emph{opposite} extreme.

\end{frame}

%----------------------------------------------------------------------
\begin{frame}{The Hypothetico-Deductive Loop}

\begin{center}
\small
\begin{tabular}{p{4.0cm}p{4.0cm}p{4.0cm}}
\toprule
\textbf{Step} & \textbf{Science} & \textbf{ML} \\
\midrule
Propose $E = h\nu$              & Theoretical hypothesis & Model architecture \\
Derive $V_0 = (h/e)\nu - W/e$   & Deduction              & Hypothesis class \\
Choose $\nu$, not $I$           & Experiment design      & Feature engineering \\
Measure $(\nu_i, V_{0i})$       & Observation            & Training data \\
Minimise RSS                    & Least squares          & MSE training \\
$|\hat\beta - h/e| < 1\%$       & Confirmation           & Test accuracy \\
Na and Li: same $\hat\beta$     & Universality test      & Zero-shot transfer \\
\bottomrule
\end{tabular}
\end{center}

\bigskip
Modern AI systems can \emph{induce} the linear law from data by
exhaustive search over a large function space.

Einstein's deductive step encoded the result of that search in one
paragraph.

\end{frame}

%----------------------------------------------------------------------
\begin{frame}{Born's Rule: An Unintended Consequence}

Einstein's quantum hypothesis required $|E|^2 \propto$ photon density.

The classical wave intensity $|E|^2$ becomes a probability density.
Einstein called this the \emph{Gespensterfeld} (ghost field).

\medskip
Max Born (1926) formalised the same logic for the Schr\"odinger wave
function:
\[
  \boxed{|\psi|^2 = \text{probability density}}
\]

\medskip
Two Nobel Prizes:
\begin{itemize}
  \item 1921: Einstein --- for the equation that made the measurement
    possible
  \item 1954: Born --- for the probabilistic interpretation the equation
    made \emph{inevitable}
\end{itemize}

\medskip
Einstein spent the rest of his life rejecting Born's rule. He won the
Nobel for a paper whose logic seeded the interpretation he most opposed.

\end{frame}

%----------------------------------------------------------------------
\begin{frame}{Summary and Conclusions}

\begin{block}{The Thesis, Revisited}
Einstein was doing machine learning in 1905---but the most important
step was the one no ML system has yet automated:
\textbf{specifying the hypothesis class}.
\end{block}

\bigskip
Three principles:
\begin{enumerate}
  \item \textbf{Theoretical priors guide design.} $E=h\nu$ fixes the
    functional form; only 5 training points suffice for 0.1\% precision.
  \item \textbf{Data quality is the binding constraint.} The regression
    is trivial; cleaning the data took 29 years.
  \item \textbf{Confirming an equation $\neq$ endorsing the model.}
    Millikan confirmed the slope to 0.5\% while rejecting the
    light-quantum interpretation.
\end{enumerate}

\bigskip
A data-driven system given only $(\nu_i, V_{0i})$ pairs can recover
$\hat\beta = 4.129\times10^{-15}$\,V$\cdot$s but cannot know it is
$h/e$. \textbf{Physical meaning requires physical theory.}

\end{frame}

%----------------------------------------------------------------------
\begin{frame}{References}

\small
\begin{itemize}
  \item Einstein, A.\ (1905). ``\"Uber einen die Erzeugung und Verwandlung
    des Lichtes betreffenden heuristischen Gesichtspunkt.''
    \textit{Ann.\ Phys.}, 17, 132--148.
  \item Lenard, P.\ (1902). ``\"Uber die lichtelektrische Wirkung.''
    \textit{Ann.\ Phys.}, 8, 149--198.
  \item Millikan, R.\,A.\ (1916). ``A Direct Photoelectric Determination
    of Planck's $h$.'' \textit{Phys.\ Rev.}, 7(3), 355--388.
  \item Born, M.\ (1926). ``Zur Quantenmechanik der
    Sto{\ss}vorg\"ange.'' \textit{Z.\ Phys.}, 37, 863--867.
  \item Hughes, A.\,L.\ (1913). ``On the Emission Velocities of
    Photo-Electrons.'' \textit{Phil.\ Trans.\ R.\ Soc.\ A}, 212, 205--226.
  \item Richardson, O.\,W.\ \& Compton, K.\,T.\ (1912). ``The
    Photoelectric Effect.'' \textit{Phil.\ Mag.}, 24, 575--594.
  \item Udrescu, S.\,M.\ \& Tegmark, M.\ (2020). ``AI Feynman.''
    \textit{Sci.\ Adv.}, 6, eaay2631.
  \item Vapnik, V.\ (1998). \textit{Statistical Learning Theory}.
    Wiley.
\end{itemize}

\end{frame}

%----------------------------------------------------------------------
\end{document}
